
% !TeX root = ../main.tex
\begin{chineseabstract}
% \footnotetext{*本研究得到某某基金（编号：）的资助}
% \astfootnote{本研究得到某某基金（编号：）的资助}
% \zhlipsum[2,5]
三维重建技术历经数十年发展已取得显著突破，并在工业检测、自动驾驶、文化遗产保护等领域广泛应用，为目标识别、场景理解等下游任务提供高精度空间信息支持。然而，传统的基于图像的三维重建方法在一些特定场景下仍存在局限性，例如物流分拣和安检领域，目标物体常被包裹覆盖，且环境光照条件受限或存在视觉遮挡，导致基于图像的方法失效。现有基于射频信号的三维重建方法多依赖复杂射频系统，硬件成本高昂。因此，开发一种兼具低成本与穿透能力的三维重建方法具有重要研究意义。本文针对上述问题，展开了如下两个方面的研究：

首先，本文提出了一种基于超宽带虚拟阵列的三维点云重建方法。该方法仅使用低成本的单发单收商用超宽带雷达对物体进行信号采集，并通过信号处理流程获取到干净有效的目标信号；针对雷达数据多样性不足的问题，设计了一种数据增强算法；通过神经网络模型从雷达信号中提取物体三维结构信息并生成多视角深度图，并通过深度图逆投影和点云融合生成重建点云。在自建数据集上的实验表明，本方法重建点云的倒角距离（Chamfer Distance, CD）误差为5.326 cm，2 cm/4 cm阈值下的F1分数分别为62.17\%与76.34\%；在完全遮挡场景中，CD误差相比于非遮挡场景仅增加0.977 cm，验证了方法对遮挡的鲁棒性。

其次，提出了一种基于仿真数据预训练的迁移学习策略。通过基于物理模拟的雷达信号仿真算法和基于深度学习的信号精炼神经网络，本文构建了大规模仿真数据集，之后采取“预训练-微调”的训练策略，利用仿真数据预训练模型参数后，在真实数据上做微调，以提升模型的数据效率、训练速度和性能表现。迁移学习训练策略下，即使仅保留真实数据集中5\%的训练数据进行微调，模型的CD误差仍能达到5.1747 cm；模型的训练周期由50个epoch缩短为仅8个epoch；同时，全量真实数据微调下，CD误差进一步降低至 4.868cm。

% \begin{enumerate}[wide,]
%     \item \zhlipsum[1]
%     \item \zhlipsum[3]
% \end{enumerate}

\chinesekeywordstype{脉冲超宽带雷达；三维重建；深度学习；迁移学习；}{应用研究}

\end{chineseabstract}

